% --- CORRECCIÓN IMPORTANTE EN LA PRIMERA LÍNEA ---
% Agregamos "xcolor=table" para que funcione \rowcolor sin errores
\documentclass[aspectratio=169, 10pt, xcolor=table]{beamer}

% --- Configuración del Tema (estilo profesional como el ejemplo) ---
\usetheme{Madrid}
\usecolortheme{dolphin}

% --- Paquetes necesarios ---
\usepackage[utf8]{inputenc}
\usepackage[spanish]{babel}
\usepackage{amsmath, amssymb}
\usepackage{booktabs}
\usepackage{graphicx}
\usepackage{listings}
\usepackage{tikz}
\usetikzlibrary{arrows.meta, positioning, shapes.geometric}

% --- Ruta de Imágenes (crea la carpeta "imagenes" y coloca ahí tus archivos) ---
\graphicspath{{./imagenes/}}

% --- Estilo para código Python ---
\definecolor{codegreen}{rgb}{0,0.6,0}
\definecolor{codegray}{rgb}{0.5,0.5,0.5}
\definecolor{codepurple}{rgb}{0.58,0,0.82}
\definecolor{backcolour}{rgb}{0.95,0.95,0.92}
\lstset{
    language=Python,
    backgroundcolor=\color{backcolour},
    commentstyle=\color{codegreen},
    keywordstyle=\color{blue},
    numberstyle=\tiny\color{codegray},
    stringstyle=\color{codepurple},
    basicstyle=\ttfamily\scriptsize,
    breaklines=true,
    frame=single,
    showstringspaces=false
}

% --- Pie de página personalizado ---
\makeatletter

\setbeamertemplate{footline}{
  \leavevmode%
  \hbox{%
  \begin{beamercolorbox}[wd=.333333\paperwidth,ht=2.25ex,dp=1ex,center]{author in head/foot}%
    \usebeamerfont{author in head/foot}\insertshortauthor
  \end{beamercolorbox}%
  \begin{beamercolorbox}[wd=.333333\paperwidth,ht=2.25ex,dp=1ex,center]{title in head/foot}%
    \usebeamerfont{title in head/foot}Sesión 11
  \end{beamercolorbox}%
  \begin{beamercolorbox}[wd=.333333\paperwidth,ht=2.25ex,dp=1ex,right]{date in head/foot}%
    \usebeamerfont{date in head/foot}NLP \hfill \insertframenumber/\inserttotalframenumber \hspace*{0.3cm}
  \end{beamercolorbox}}%
  \vskip0pt%
}

\makeatother

% --- Datos de la portada ---
\title[SESIÓN 11 NLP]{\textbf{Sesión 11}}
\subtitle{Tareas Avanzadas de NLP: NER, Question Answering y Summarization}
\author[Prof. C. Sánchez]{Carlos César Sánchez Coronel}
\institute{\textbf{NLP}}
\date{}

% Logo opcional (descomentar si tienes la imagen)
% \titlegraphic{\includegraphics[height=1.5cm]{logo.png}}

% --- Eliminar la barra de navegación (opcional) ---
\setbeamertemplate{navigation symbols}{}

% --- Índice al inicio de cada sección ---
\AtBeginSection[]{
  \begin{frame}
    \frametitle{Contenido}
    \tableofcontents[currentsection]
  \end{frame}
}

% ============================================================================
% INICIO DEL DOCUMENTO
% ============================================================================
\begin{document}

% --- Portada ---
\begin{frame}
    \titlepage
\end{frame}

% --- Índice general ---
\begin{frame}{Contenido}
    \tableofcontents
\end{frame}

% ============================================================================
% SECCIÓN 1: NAMED ENTITY RECOGNITION (NER)
% ============================================================================
\section{Named Entity Recognition (NER)}

\begin{frame}{Definición y objetivos}
    \begin{itemize}
        \item Identificar y clasificar menciones de entidades con nombre en texto.
        \item Entidades típicas: personas (PER), organizaciones (ORG), lugares (LOC), fechas (DATE), cantidades (MONEY), etc.
        \item Aplicaciones: búsqueda semántica, extracción de información, análisis de redes sociales.
    \end{itemize}
    % Basado en Pregunta 1
\end{frame}

\begin{frame}{Sistemas de etiquetado: BIO}
    \begin{block}{Esquema BIO}
        \begin{itemize}
            \item \textbf{B} (Beginning): inicio de una entidad.
            \item \textbf{I} (Inside): dentro de una entidad.
            \item \textbf{O} (Outside): fuera de entidad.
        \end{itemize}
    \end{block}
    \begin{exampleblock}{Ejemplo}
        ``Banco Santander abrió oficina en Madrid'' \\
        \vspace{0.2cm}
        Banco (B-ORG) Santander (I-ORG) abrió (O) oficina (O) en (O) Madrid (B-LOC)
    \end{exampleblock}
    % Basado en Pregunta 1
\end{frame}

\begin{frame}{Arquitectura típica con BERT}
    \begin{itemize}
        \item Modelo pre-entrenado (BERT) más una capa lineal de clasificación por token.
        \item Entrada: tokens con subpalabras.
        \item Salida: logits para cada token, softmax por token.
        \item Pérdida: entropía cruzada por token (ignorando subpalabras no deseadas).
    \end{itemize}
    % Basado en Pregunta 2
    % IMAGEN: Diagrama de NER con BERT
\end{frame}

\begin{frame}{Métrica de evaluación: F1 a nivel de entidad}
    \begin{itemize}
        \item No se usa accuracy por token (puede ser engañoso por muchas etiquetas ``O'').
        \item Se considera una entidad correcta si el tipo y los límites coinciden exactamente con la referencia.
        \item Se calculan verdaderos positivos (VP), falsos positivos (FP) y falsos negativos (FN) sobre entidades completas.
        \item Precisión = VP/(VP+FP), Recall = VP/(VP+FN), F1 = media armónica.
    \end{itemize}
    % Basado en Pregunta 3
\end{frame}

\begin{frame}[fragile]{Ejemplo con Hugging Face}
    \begin{lstlisting}
from transformers import pipeline

ner_pipeline = pipeline("ner", model="dslim/bert-base-NER")
text = "Steve Jobs founded Apple in Cupertino."
entities = ner_pipeline(text)
for entity in entities:
    print(entity)
    \end{lstlisting}
    % Basado en Pregunta 15
\end{frame}

% ============================================================================
% SECCIÓN 2: QUESTION ANSWERING (QA)
% ============================================================================
\section{Question Answering (QA)}

\begin{frame}{Tipos de QA}
    \begin{columns}[t]
        \column{0.5\textwidth}
        \textbf{Extractivo}
        \begin{itemize}
            \item La respuesta es un fragmento del contexto.
            \item Ejemplo: SQuAD.
            \item Ventaja: fiable, no inventa.
            \item Desventaja: limitado a texto presente.
        \end{itemize}
        \column{0.5\textwidth}
        \textbf{Generativo}
        \begin{itemize}
            \item La respuesta se genera libremente.
            \item Ejemplo: sistemas de diálogo.
            \item Ventaja: flexible, puede parafrasear.
            \item Desventaja: riesgo de alucinaciones.
        \end{itemize}
    \end{columns}
    % Basado en Pregunta 5
\end{frame}

\begin{frame}{QA extractivo con BERT}
    \begin{itemize}
        \item Entrada: ``[CLS] pregunta [SEP] contexto [SEP]''.
        \item Dos cabezas lineales adicionales:
        \begin{itemize}
            \item Una predice logits de inicio.
            \item Otra predice logits de fin.
        \end{itemize}
        \item Softmax sobre todos los tokens para inicio y fin.
        \item Pérdida: suma de entropías cruzadas de inicio y fin.
        \item Inferencia: seleccionar span $(i,j)$ que maximice $p_{\text{start}}(i) \cdot p_{\text{end}}(j)$ con $i \le j$.
    \end{itemize}
    % Basado en Pregunta 6
\end{frame}

\begin{frame}{Métricas para QA extractivo}
    \begin{block}{Exact Match (EM)}
        Porcentaje de respuestas que coinciden exactamente (carácter por carácter) con la referencia.
    \end{block}
    \begin{block}{F1-score}
        Media armónica de precisión y recall sobre tokens de la respuesta.
    \end{block}
    \begin{exampleblock}{Ejemplo}
        Referencia: ``Steve Jobs'', Predicción: ``Steve'' \\
        EM = 0, F1 = 0.667 (precisión 1, recall 0.5)
    \end{exampleblock}
    % Basado en Pregunta 7
\end{frame}

\begin{frame}[fragile]{Ejemplo con Hugging Face}
    \begin{lstlisting}
from transformers import pipeline

qa_pipeline = pipeline("question-answering", model="distilbert-base-cased-distilled-squad")
context = "Steve Jobs founded Apple in Cupertino."
question = "Who founded Apple?"
result = qa_pipeline(question=question, context=context)
print(f"Answer: {result['answer']}, score: {result['score']}")
    \end{lstlisting}
    % Basado en Pregunta 15
\end{frame}

% ============================================================================
% SECCIÓN 3: SUMMARIZATION
% ============================================================================
\section{Summarization}

\begin{frame}{Tipos de resumen}
    \begin{columns}[t]
        \column{0.5\textwidth}
        \textbf{Extractivo}
        \begin{itemize}
            \item Selecciona oraciones del texto original.
            \item Más simple, garantiza fluidez.
            \item Puede ser menos cohesivo.
        \end{itemize}
        \column{0.5\textwidth}
        \textbf{Abstractivo}
        \begin{itemize}
            \item Genera nuevas oraciones parafraseando.
            \item Más flexible, puede condensar mejor.
            \item Riesgo de alucinaciones.
        \end{itemize}
    \end{columns}
    % Basado en Pregunta 9
\end{frame}

\begin{frame}{Arquitectura para summarization abstractivo}
    \begin{itemize}
        \item Modelos seq2seq pre-entrenados: BART, T5.
        \item Encoder: bidireccional (procesa documento).
        \item Decoder: autoregresivo (genera resumen).
        \item Atención cruzada entre decoder y encoder.
    \end{itemize}
    % Basado en Pregunta 10
    % IMAGEN: Diagrama encoder-decoder
\end{frame}

\begin{frame}{Métrica: ROUGE}
    \begin{block}{ROUGE-N}
        Solapamiento de n-gramas entre resumen generado y referencia.
    \end{block}
    \begin{block}{ROUGE-L}
        Basado en la subsecuencia común más larga (LCS).
    \end{block}
    \begin{block}{ROUGE-S}
        Skip-bigrams (pares de palabras con gap).
    \end{block}
    % Basado en Pregunta 11
\end{frame}

\begin{frame}{Ejemplo numérico de ROUGE}
    \textbf{Referencia}: ``El gato persigue al ratón'' \\
    \textbf{Generado}: ``El gato caza ratones'' \\
    \begin{itemize}
        \item ROUGE-1 (recall): unigramas comunes: El, gato (2/5 = 0.4)
        \item ROUGE-2 (recall): bigramas comunes: ``El gato'' (1/4 = 0.25)
        \item ROUGE-L: LCS = ``El gato'' (2 palabras) → recall = 2/5 = 0.4
    \end{itemize}
    % Basado en Pregunta 11
\end{frame}

\begin{frame}[fragile]{Ejemplo con Hugging Face}
    \begin{lstlisting}
from transformers import pipeline

summarizer = pipeline("summarization", model="facebook/bart-large-cnn")
text = "Elon Musk, CEO of Tesla and SpaceX, announced that Tesla's new factory in Berlin will begin production next month. The factory, which cost over $5 billion, is expected to create 10,000 jobs."
summary = summarizer(text, max_length=50, min_length=25)
print(summary[0]['summary_text'])
    \end{lstlisting}
    % Basado en Pregunta 15
\end{frame}

% ============================================================================
% SECCIÓN 4: PREPARACIÓN PARA ENTREVISTAS
% ============================================================================
\section{Preparación para Entrevistas}

\begin{frame}{Preguntas típicas}
    \begin{itemize}
        \item ¿Qué métrica usarías para evaluar un sistema de NER? (F1 a nivel de entidad)
        \item Explica la diferencia entre QA extractivo y generativo.
        \item ¿Qué es ROUGE y cómo se calcula?
        \item ¿Qué problemas tiene el summarization abstractivo? (alucinaciones, repeticiones)
    \end{itemize}
    % Basado en PDF
\end{frame}

\begin{frame}{Preguntas avanzadas}
    \begin{itemize}
        \item ¿Cómo manejarías contextos largos en QA con modelos de ventana limitada? (división, Longformer)
        \item ¿Qué es el esquema BILOU y cómo mejora a BIO?
        \item ¿Cómo detectarías alucinaciones en summarization? (métricas de factualidad)
        \item ¿Qué modelos pre-entrenados existen para QA sobre tablas? (TAPAS)
    \end{itemize}
    % Basado en Preguntas 8, 12, 13, 14, 18
\end{frame}

% ============================================================================
% SECCIÓN 5: CONEXIÓN CON EL RESTO DEL CURSO
% ============================================================================


% ============================================================================
% SECCIÓN 6: RESUMEN
% ============================================================================
\section{Resumen}

\begin{frame}{Lo que hemos aprendido}
    \begin{exampleblock}{Conceptos clave}
        \begin{itemize}
            \item \textbf{NER}: identificación de entidades con esquemas BIO/BILOU; arquitectura con BERT; evaluación con F1 a nivel de entidad.
            \item \textbf{QA extractivo}: BERT con dos cabezas para inicio y fin; métricas EM y F1.
            \item \textbf{Summarization abstractivo}: modelos seq2seq (BART, T5); métrica ROUGE.
            \item \textbf{Implementación}: pipelines de Hugging Face para cada tarea.
            \item \textbf{Desafíos}: alucinaciones, contextos largos, datos desbalanceados.
        \end{itemize}
    \end{exampleblock}
\end{frame}

% --- Diapositiva final ---
\begin{frame}
    \centering
    \Huge \textbf{¡Gracias!}

\end{frame}

\end{document}